\subsection*{Notation, assumptions, and the two anchoring conventions}

This subsection is the canonical register for the derivation record D0 through
D9, with D3b the adaptive companion of D3: D0 is the model and cost lemma, D1
information saturation, D2 the exchange rate, D3 the minimax lower bounds, D4
and D5 the design geometry, D6 safe certainty equivalence, D7 the anchoring
crossover, and D8 and D9 the ROI analysis and the instantiation. Every symbol
is defined once, here; every assumption is numbered here and is cited by
number in the derivation sections and in the register statements (T1, L2,
\dots). Where a later section restates a condition, this register is
authoritative. The formal proofs of the register statements are collected in
the proofs appendix.

\paragraph{Notation.}
The symbol table is reproduced in three blocks: the retraining loop and its
objective, the stochastic and informational quantities, and the symbols
specific to the design and economic analyses. The third column records where
each symbol is introduced: a derivation section (D0 through D9), the base
structural market theory (cited as REFLEX \citep{reflex2026} with its section numbers,
a third-party source), or an entry of the assumption register below (the
trust-region radius $r$ is fixed by A4). Two register conventions carry over
unchanged: all logarithms are natural, and a prime denotes the derivative of
a scalar function (in the ASCII sources it also denotes the matrix transpose;
in this appendix transposes are written $(\cdot)^\top$, so the prime is
unambiguous).

\begin{center}\small
\begin{tabular}{@{}l>{\raggedright\arraybackslash}p{4.9cm}l@{}}
\hline
Symbol & Meaning & Defined in / source\\
\hline
$h$, $h_t$ & deployed decision (half-spread); its value at
  step $t$ & D0\\
$\hstar$ & performatively stable point (RRM fixed point)
  & REFLEX \citep{reflex2026}\\
$\hsp$, $\hpo$ & stable point ($\hsp = \hstar$); performative
  optimum & REFLEX \citep{reflex2026}\\
$\dev_t$ & deviation $\dev_t = h_t - \hstar$ & D0\\
$m$ & retraining modulus $m = \eps\beta/\gamma \in [0,1)$
  & REFLEX \citep{reflex2026}\\
$\gamma$, $\beta$, $\eps$, $\psi$ & curvature, smoothness,
  response slope, adverse severity & REFLEX \citep{reflex2026} (closed forms)\\
$\Phi$ & true performative objective $\Phi(h) = J(h;\tau(h))$
  & REFLEX \citep{reflex2026}\\
$\delta_0$ & $\Phi'(\hstar) = -\beta(\hstar-\psi)\,\eps(\hstar)$,
  nonzero & D0 Sec.~1\\
$\gpo$ & $-\Phi''$ at the operating point & REFLEX \citep{reflex2026}\\
$M$ & $d$-dimensional loop Jacobian, $\rho(M) < 1$ & D1 Sec.~2\\
$P$, $P_u$ & discrete-Lyapunov energy matrices & D1 Sec.~2\\
$\Gpo$ & $d$-dimensional objective curvature matrix & D9\\
\hline
\end{tabular}
\end{center}

\begin{center}\small
\begin{tabular}{@{}l>{\raggedright\arraybackslash}p{4.9cm}l@{}}
\hline
Symbol & Meaning & Defined in / source\\
\hline
$\sigma$ & response-noise standard deviation ($\sigma_\tau$ in
  REFLEX) & D0\\
$\sigma_e$, $u_t$ & exploration intensity; exploration input & D0\\
$\xi_t$, $\zeta_t$ & deployment and response noise (iid, mean
  zero) & D0\\
$\phi$ & noise-feedback gain (observed noise flowing into the
  next deployment) & D2 Sec.~4\\
$v$ & stationary deviation variance $\sigma_e^2/(1-m^2)$ & D2\\
$v_\phi$ & the same with feedback:
  $(\sigma_e^2 + \phi^2\sigma^2)/(1-m^2)$ & D2 Sec.~4\\
$\Sxx$ & centered design energy
  $\sum_t (\dev_t - \bar{\dev})^2$ & D2\\
$D_T$ & deviation budget $\sum_t \dev_t^2$ & D0 Sec.~2\\
$C_T$ & value budget: expected incremental cost (L1)
  & D0 Secs.~2--3\\
$I_T$ & Fisher information for $\eps$ & D1/D3\\
$\mathcal{F}_t$ & history sigma-field
  $\sigma(\dev_0, u_{<t}, \tau_{<t})$ & D0\\
$\delta$, $s$ & two-point prior half-width; its sign (minimax)
  & D3\\
$g_1$ & $\beta\,\lvert \hstar - \psi \rvert$, coefficient of the
  unknown linear part & D3\\
$\TV$, $\KL$ & total variation; Kullback--Leibler divergence & D3\\
$r$ & trust-region radius on $\lvert \dev_t \rvert$ & A4\\
\hline
\end{tabular}
\end{center}

\begin{center}\small
\begin{tabular}{@{}l>{\raggedright\arraybackslash}p{4.9cm}l@{}}
\hline
Symbol & Meaning & Defined in / source\\
\hline
$M$ (design) & exploration second moment $\E[\dev\,\dev^\top]$;
  appears only in D4 & D4\\
$B$ & value budget in the design problems & D4\\
$F$ & curvature dispersion
  $d\,\tr(\Gpo)/(\tr \Gpo^{1/2})^2$ & D4 (C5.1)\\
$c$ & correction-direction vector ($c$-optimality) & D4 Sec.~3\\
$s(h)$ & structural-family sensitivity vector, $p = 3$ & D5\\
$w$ & probe half-width (crossover analysis) & D7\\
$\delta_{\mathrm{mis}}$ & $C^1$ misspecification of the
  structural family & D6/D7\\
$\kappa$ & correction sensitivity
  $\lvert\, d\hpo/d\eps \,\rvert$ & D8\\
$\rho_{\mathrm{disc}}$ & discount rate (ROI analysis) & D8\\
\hline
\end{tabular}
\end{center}

Three deliberate symbol reuses deserve a caveat, recorded here so that no
derivation needs to repeat it. First, $M$ names both the loop Jacobian and
the exploration second moment; the design $M$ occurs only in D4/D5, and
context disambiguates the two. Second, $\kappa$ is the correction
sensitivity $\lvert\, d\hpo/d\eps \,\rvert$ of D8/D9, not a multi-dealer
spillover parameter: no such symbol occurs in this paper. Third, $\phi$ is
the noise-feedback gain of D2, while the T9 material in D3 separately
writes $\phi_u$ and $\phi_0$ for elements of the sensitivity span, a
self-contained reuse local to that subsection.

\paragraph{Assumption register.}
Assumptions are cited by number (A1 through A6) throughout the record.

\begin{itemize}
\item \textbf{A1 (local quadratic scope).} $\Phi$ is $C^3$ and strongly
concave on the operating interval; all expansions are taken at $\hstar$, with
the third-order remainder controlled through $L_3 = \sup \lvert \Phi'''
\rvert$. The identities of the record are exact for the quadratic model and
first-order accurate otherwise. As a matter of method, the nonlinear
deviation is measured (the drift study), never assumed away.
\item \textbf{A2 (exploration classes).} All policies are non-anticipating:
$u_t$ is $\mathcal{F}_t$-measurable. Two subclasses are used.
\emph{A2-sym}, conditionally symmetric exploration,
$\E[u_t \mid \mathcal{F}_t] = 0$, under which the linear term of L1 vanishes
in expectation. \emph{A2-adaptive}, unrestricted non-anticipating
exploration, under which costs are measured from the certainty-equivalent
anchor (D0 Sec.~4; convention 2 below) and the exploitation lemma L2 of D3b
applies.
\item \textbf{A3 (noise exogeneity, with its named relaxation).} The
response noise $\zeta_t$ is iid, mean zero, with variance $\sigma^2$,
independent of $\mathcal{F}_t$. Relaxation A3$'$: the feedback channel
$\phi\,\zeta_t$ enters $\dev_{t+1}$. Under A3$'$ the exact $O(1/T)$ bias of
D2 Sec.~4 applies, with its $(3m-1)$ sign structure (P2.2), and
``exploration-dominant design'' means $\sigma_e^2 \gg \phi^2 \sigma^2$.
\item \textbf{A4 (trust region).} $\lvert \dev_t \rvert \le r$ pathwise.
The radius $r$ enters the constants of D3b (L2, T4), the feasible set of
the design problems (T5), and the basin argument of D6.
\item \textbf{A5 (stationary-scale exploration; D3b only).}
$S_T = \Theta(T)$, with $S_T = \sum_{t \le T} \dev_t^2$ the cumulative
design energy. The transient-dominated regime $S_T = o(T)$ is governed by
D1 (T1) instead.
\item \textbf{A6 (stable base loop).} $m < 1$ in the scalar case and
$\rho(M) < 1$ in the $d$-dimensional case, for all stationary-regime
statements. Boundary behaviour ($m = 1$) appears only in the degeneracy
statement C1.2.
\end{itemize}

\paragraph{The two anchoring conventions (audit-critical).}
$C_T$ is an incremental cost, and its value depends on the baseline against
which the increment is measured. The record fixes one anchor per exploration
class; these are the two audit-critical definitions.

\emph{Convention 1 (incremental anchor; all A2-sym results).} $C_T$ is
measured relative to the jitter-free loop at $\hstar$, that is, the same
retraining loop run without exploration. Anchoring at $\hpo$ instead adds
the sunk echo-chamber term and provably destroys the invariance of D2 (T2):
with anchor gap $g = \hstar - \hpo$, the anchored exchange-rate product is
inflated by exactly the factor
\begin{equation*}
1 + \frac{g^2}{v},
\end{equation*}
a computation verified in D0. The formal proof is recorded in the proofs
appendix.

\emph{Convention 2 (certainty-equivalent anchor; all A2-adaptive results).}
$C_T$ is measured relative to the deployment that is optimal under current
posterior knowledge. Exploitation of learned structure is thereby
re-anchored away, and D3b (L2, T4) prices only the unlearned residual.

The two anchors coincide under A2-sym with an uninformative start, so the
conventions agree on the class where both are defined.

\subsection*{D0: The model, the two budgets, and the cost-equivalence lemma (L1)}

This subsection records the local model, the two exploration budgets, and the
computation behind the cost-equivalence lemma L1. It is derived in full and
verified numerically (\texttt{verify\_all.py}, section V0). Everything
downstream (D1 through D8) consumes the definitions and the lemma established
here. Notation follows the REFLEX market model \citep{reflex2026}
($\gamma$, $\beta$, $\eps$, $\psi$, $\hsp$, $\hpo$).

\paragraph{The local model.}
Work in a neighbourhood of the performatively stable point $\hstar$, the RRM
fixed point of that market \citep{reflex2026}, and write deviations
$\dev_t := h_t - \hstar$. The linearized closed loop is
\begin{align*}
  \dev_{t+1} &= -m\,\dev_t + u_t
    && \text{(retraining cobweb plus exploration)},\\
  \tau_t &= \tau(\hstar) - \eps\,\dev_t + \zeta_t
    && \text{(observed flow)}.
\end{align*}
Here $m = \eps\beta/\gamma \in [0,1)$ is the retraining modulus, whose
closed form comes with the model \citep{reflex2026}. The exploration input $u_t$ is chosen by the operator:
it may be iid noise ($u_t = \sigma_e\,\xi_t$), a deterministic schedule, or
adaptive, that is $\mathcal{F}_t$-measurable with
$\mathcal{F}_t = \sigma(\dev_0, u_0, \dots, u_{t-1}, \tau_0, \dots,
\tau_{t-1})$. The response noise $\zeta_t$ is iid, mean zero, with variance
$\sigma^2$, independent of $\mathcal{F}_t$; the feedback violation of this
independence assumption is quantified in the exchange-rate derivation D2
(the feedback-bias result P2.2). The estimand is the response
slope $\eps$, scalar here; the $d$-dimensional version (D4, behind T5) and
the structural-family version (D5, behind T6) are treated in their own
derivation subsections.

The true performative objective $\Phi(h) = J(h; \tau(h))$ (REFLEX theory
1.2) is smooth and strongly concave on the operating interval, with
\begin{equation*}
  \Phi'(\hstar) = \Delta(\hstar) = -\beta\,(\hstar - \psi)\,\eps(\hstar)
  =: \delta_0, \qquad
  \Phi''(\hstar) = -\gpo
\end{equation*}
(a direct computation in the market model of \citet{reflex2026}). The
slope $\delta_0$ is nonzero precisely
because $\hstar \neq \hpo$.

\paragraph{The two budgets.}
Two notions of exploration budget are used throughout:
\begin{align*}
  D_T &:= \sum_{t\le T} \dev_t^2
    && \text{(deviation budget: realized design energy)},\\
  C_T &:= \E\Big[\sum_{t\le T} \big(\Phi(\hstar) - \Phi(h_t)\big)\Big]
    && \text{(value budget: expected incremental cost)}.
\end{align*}
$D_T$ is the tracking-error form, a desk's deviation limit: it is
observable, its meaning is policy-independent, and it is the correct
constraint for the exact minimax theorem T3 (derivation D3). $C_T$ is the
realized-P\&L form. Lemma L1 relates the two.

\textbf{Lemma L1 (cost equivalence).} For any exploration rule that is
non-anticipating and conditionally symmetric, meaning
$\E[u_t \mid \mathcal{F}_t] = 0$ (for example iid exploration, or a
scheduled sign-symmetric jitter), in the regime where third-order terms are
negligible,
\begin{equation*}
  C_T = \tfrac{1}{2}\,\gpo\,\E[D_T] + R_T, \qquad
  |R_T| \le \frac{L_3}{6}\,\E\Big[\sum_{t\le T} |\dev_t|^3\Big],
\end{equation*}
with $L_3 = \sup |\Phi'''|$ on the operating interval. The linear term of
the expansion contributes exactly zero to the expectation, and contributes
a fluctuation of standard deviation
$|\delta_0|\,\sqrt{\Var\big(\sum_{t\le T}\dev_t\big)}$ to the realized
cost.

\textbf{Proof.} Taylor-expand $\Phi$ at $\hstar$:
\begin{equation*}
  \Phi(\hstar) - \Phi(h_t)
  = -\delta_0\,\dev_t + \tfrac{1}{2}\,\gpo\,\dev_t^2 + r_t,
  \qquad |r_t| \le \tfrac{L_3}{6}\,|\dev_t|^3,
\end{equation*}
then sum over $t \le T$ and take expectations. For the linear term: $\dev_t$
is built from $\{u_s\}_{s<t}$ through the linear recursion, so the mean
obeys $\E[\dev_{t+1}] = -m\,\E[\dev_t] + \E[u_t]$; conditional symmetry
gives $\E[u_t] = \E\big[\E[u_t \mid \mathcal{F}_t]\big] = 0$, hence
$\E[\dev_t] = (-m)^t\,\E[\dev_0]$. Starting at the operating point
($\dev_0 = 0$), or in the stationary regime, $\E[\dev_t] = 0$ for all $t$,
so $\E\big[\sum_{t\le T} \delta_0\,\dev_t\big] = 0$. The quadratic term
yields $\tfrac{1}{2}\gpo\,\E[D_T]$, and the remainder is bounded by the
stated third-order term. \qedsymbol

\paragraph{Fluctuation decomposition (exact, and a corrected earlier
claim).}
The realized cost's variance splits into two terms with zero cross term:
the cross-moment sums built from $\E[\dev_t\,\dev_s^2]$ vanish by the sign
symmetry of the stationary law. Exactly,
\begin{equation*}
  \Var\big(C_T^{\mathrm{real}}\big)
  = \delta_0^2\,\Var\Big(\sum_{t\le T}\dev_t\Big)
  + \frac{\gpo^2}{4}\,\Var\Big(\sum_{t\le T}\dev_t^2\Big),
\end{equation*}
with the long-run scales, $v$ denoting the stationary variance of $\dev_t$,
\begin{align*}
  \Var\Big(\sum_{t\le T}\dev_t\Big)
    &\sim T\,v\,\frac{1-m}{1+m}
    && \text{(long-run variance: negative autocorrelation helps)},\\
  \Var\Big(\sum_{t\le T}\dev_t^2\Big)
    &\sim T\,\frac{2\,v^2\,(1+m^2)}{1-m^2}
    && \text{(Gaussian stationary AR(1))}.
\end{align*}
A caveat, recorded per the falsification convention: an earlier draft
asserted that the linear term always dominates. The first verification run
falsified that at moderate $\delta_0$ (measured sd ratio $1.62$, exactly
reproduced by the two-term formula above), and the claim is now stated as
the full decomposition. Which term dominates is decided by comparing
$\delta_0^2\,(1-m)/(1+m)$ against
$(\gpo^2\,v/2)\,(1+m^2)/(1-m^2)$, both computable. V0 now verifies the
decomposition itself.

\paragraph{Why the anchoring is load-bearing (recorded from the audit).}
Anchoring the cost at the performative optimum $\hpo$ instead of $\hstar$
adds terms of the type $T\,(\hstar - \hpo)^2$ that are paid whether or not
the operator explores: the sunk echo-chamber cost of blindness (REFLEX
theory 1.2, 6b). Numerically (\texttt{verify\_m2\_identity.py}, Claim 3):
writing $g = \hstar - \hpo$ for the gap, the mis-anchored product equals
\begin{equation*}
  \tfrac{1}{2}\,\gpo\,\sigma^2\,\Big(1 + \frac{g^2}{v}\Big),
\end{equation*}
which destroys the exchange-rate invariance of D2 (theorem T2). The
incremental anchor at $\hstar$ is the definition under which every later
theorem is true.

\paragraph{The certainty-equivalent anchor (needed for the adaptive
minimax bound T4, derivation D3).}
For adaptive policies the symmetric-exploration condition can be violated
deliberately: since $\delta_0 \neq 0$, drifting toward $\hpo$ earns
first-order value. Two observations discipline this.
First, decompose $\Phi'(\hstar)$ as a function of the estimand:
\begin{equation*}
  \varphi(\eps) := \Phi'(\hstar; \eps) = -\beta\,(\hstar - \psi)\,\eps,
\end{equation*}
which is linear in $\eps$. Under a prior centred at $\bar\eps$, split
$\varphi(\eps) = \varphi(\bar\eps)
+ \varphi'(\bar\eps)\,(\eps - \bar\eps)$.
Second, the first (known) part is exploitable without any learning: it is a
pure consequence of already-known information, that is, of choosing the
anchor wrongly. Define therefore the certainty-equivalent anchor
$\hstar_{\mathrm{CE}}$: the deployment that is optimal under current
knowledge (the posterior mean), with exploration cost measured relative to
the trajectory that sits at $\hstar_{\mathrm{CE}}$. Relative to that anchor
the residual linear coefficient is
\begin{equation*}
  \varphi'(\bar\eps)\,(\eps - \bar\eps)
  = -\beta\,(\hstar - \psi)\,(\eps - \bar\eps),
\end{equation*}
which has zero prior mean: only genuinely unknown structure can be
exploited, and the exploitation--information lemma L2 (in D3, feeding T4)
bounds that exploitation by the information already purchased. This
definition also matches practice: a desk's exploration cost is measured
against its own current best policy, not against an oracle's.

\paragraph{Numerical verification (V0).}
Three checks close the record for this subsection.
\begin{itemize}
  \item $\E[C_T] = \tfrac{1}{2}\,\gpo\,\E[D_T]$ holds within Monte Carlo
    error, for a quadratic-plus-linear objective with $\delta_0 \neq 0$,
    under iid exploration.
  \item The realized-cost fluctuation scale matches
    $|\delta_0|\,\mathrm{sd}\big(\sum_{t\le T}\dev_t\big)$: the linear
    term is a variance contribution, not a bias.
  \item A deliberately asymmetric (drifting) exploration rule breaks the
    lemma in the predicted direction and by the predicted first-order
    amount $-\delta_0\,\E\big[\sum_{t\le T}\dev_t\big]$: the
    conditional-symmetry condition of L1 is necessary, not decorative.
\end{itemize}

\subsection*{D1: Information saturation by the Lyapunov identity (T1, C1.1, C1.2, R1)}

This subsection records the derivation behind T1 and its corollaries C1.1 and
C1.2, together with the non-normal refinement R1. Status of the record: the
scalar closed form is verified to $10^{-16}$, and the $d$-dimensional
identity, including the non-normal amplification, is verified numerically
(V1); the proofs are complete.

\paragraph{Scalar warm-up (the verified base case).}
With no exploration ($u_t = 0$) and modulus $m < 1$, the deviation obeys
$\dev_t = (-m)^t \dev_0$, so the design energy accumulated by horizon $T$ is
\begin{equation*}
D_T \;=\; \sum_{t=0}^{T} \dev_t^2
\;=\; \dev_0^2\,\frac{1 - m^{2(T+1)}}{1 - m^2}
\;\longrightarrow\; \frac{\dev_0^2}{1 - m^2}
\qquad (T \to \infty).
\end{equation*}
The Fisher information for the response slope $\eps$ from OLS on the
trajectory is $\Sxx/\sigma^2 \le D_\infty/\sigma^2$: bounded uniformly in
the horizon. Hence, by Cram\'er--Rao, for any unbiased estimator and at
every horizon,
\begin{equation*}
\Var(\hat\eps) \;\ge\; \frac{\sigma^2\,(1 - m^2)}{\dev_0^2}.
\end{equation*}
This is the scalar case of T1; with $\dev_0 = h_0 - \hstar$ the cap
$\dev_0^2/(1-m^2)$ is the quantity in the register statement of C1.1.

\paragraph{C1.1 (safety implies blindness).}
The cap $\dev_0^2/(1-m^2)$ is strictly increasing in $m$: the
fast-contracting (safe) loop reveals least about its own performativity, and
identifiability is a near-instability phenomenon. This is immediate from the
closed form above; the economic content is stated in the paper body.

\paragraph{C1.2 (design degeneracy).}
At $m = 1$ the noiseless iterates form the period-two orbit
$\{+\dev_0, -\dev_0\}$; for $m < 1$ the support collapses geometrically onto
$\{0\}$. In either case the empirical design measure has at most two
effective support points, which by the Chebyshev counting argument of D5
(theorem T6) cannot identify a three-parameter response family. Curvature is
invisible to retraining trajectories at every $m$.

\paragraph{The $d$-dimensional identity via the discrete Lyapunov equation.}
Let the linearized joint loop be $\dev_{t+1} = M \dev_t$ with
$M \in \mathbb{R}^{d \times d}$ and $\rho(M) < 1$; in the multi-bond and
multi-dealer instantiations $M$ comes from the REFLEX market model
\citep{reflex2026} and is generally not symmetric.

\textbf{Theorem (exact energy; T1).} The total design energy is the
quadratic form
\begin{equation*}
E(\dev_0) \;:=\; \sum_{t \ge 0} \big\| M^t \dev_0 \big\|^2
\;=\; \dev_0^\top P\,\dev_0,
\end{equation*}
where $P$ is the unique positive semidefinite solution of the discrete
Lyapunov equation
\begin{equation*}
P \;=\; I + M^\top P M.
\end{equation*}

\textbf{Proof.} Set $P := \sum_{t \ge 0} (M^\top)^t M^t$; the series
converges absolutely if and only if $\rho(M) < 1$ (Gelfand's spectral
radius formula). Then
\begin{equation*}
M^\top P M \;=\; \sum_{t \ge 0} (M^\top)^{t+1} M^{t+1} \;=\; P - I,
\end{equation*}
so $P$ solves the equation. Uniqueness: the difference $Q$ of two solutions
satisfies $Q = M^\top Q M = (M^\top)^t Q\, M^t \to 0$. Finally,
$\dev_0^\top P\,\dev_0 = \sum_{t \ge 0} \|M^t \dev_0\|^2$ by expanding the
sum term by term. \qedsymbol

Scalar check: $P = 1/(1 - m^2)$, recovering the warm-up above.

\paragraph{Directional information.}
For the response slope in a direction $u$ (the observation
$\tau_t = \tau^{*} - \eps\, u^\top \dev_t + \zeta_t$ probes the component
along $u$), the relevant energy is
\begin{equation*}
\sum_{t \ge 0} \big( u^\top M^t \dev_0 \big)^2
\;=\; \dev_0^\top P_u\,\dev_0,
\qquad
P_u \;=\; u u^\top + M^\top P_u M,
\end{equation*}
the same Lyapunov structure with a rank-one right-hand side. By
Cram\'er--Rao this yields
$\Var(\hat\eps) \ge \sigma^2 \big( \dev_0^\top P_u\,\dev_0 \big)^{-1}$ for
the component along any direction $u$, which is the bound stated in T1.
Hence the Fisher information is bounded in every direction, and the
information operator satisfies
$\mathcal{I} \preceq P_{\mathrm{full}} / \sigma^2$, where
$P_{\mathrm{full}}$ solves the matrix-valued Lyapunov equation. Saturation
is not a scalar accident.

\paragraph{Non-normality: transiently amplified information (R1).}
If $M$ is normal, $P$ has eigenvalues $1/(1 - \lambda_i^2)$ and the energy
is governed by the spectrum alone. If $M$ is non-normal (heterogeneous
multi-dealer Jacobians are), transient growth inflates the energy beyond the
spectral prediction:
\begin{equation*}
E(\dev_0) \;\le\; \kappa(V)^2\,\frac{\|\dev_0\|^2}{1 - \rho(M)^2},
\end{equation*}
where $V$ is the eigenbasis of $M$ and $\kappa(V)$ its condition number, and
the inequality is achieved in order: for the Jordan-type family
\begin{equation*}
M \;=\; \begin{pmatrix} m & a \\ 0 & m \end{pmatrix},
\end{equation*}
the energy along the first coordinate grows like $a^2/(1 - m^2)^3$ for
large $a$, verified numerically in V1 against the exact Lyapunov solve.
Badly-conditioned ecosystems thus leak extra information about
themselves during transients, the free-information remark of the paper, now
with an exact constant ($P$ itself). The paper states the exact $P$-form and
keeps the pseudospectral refinement for the journal version.

\paragraph{Scope: what T1 does not say.}
Two caveats are part of the record and are stated in the paper.
\begin{itemize}
\item With exploration on, saturation is replaced by the exchange rate of
D2 (T2): the transient information $\dev_0^\top P\,\dev_0 / \sigma^2$ is
the free part; everything after is bought.
\item The bound is for the linearized loop; the simulator's nonlinear
regimes are handled empirically (the drift study, under the D2/V2
protocol).
\end{itemize}

\paragraph{Numerical verification (V1).}
Four checks anchor the derivation. (i) The scalar closed form against the
direct sum: exact to floating point. (ii) Random stable $M$ of sizes
$2 \times 2$ and $5 \times 5$: the direct energy sum against the Lyapunov
solve (a \texttt{scipy}-free fixed-point iteration), agreement to
$10^{-12}$. (iii) Non-normal amplification: the energy for the family
$M = \left(\begin{smallmatrix} m & a \\ 0 & m \end{smallmatrix}\right)$
grows as $a^2$, far exceeding the normal-matrix bound $1/(1 - m^2)$, and
matches $\dev_0^\top P\,\dev_0$ exactly. (iv) The directional version: the
rank-one Lyapunov solution $P_u$ against the direct directional sum.

\subsection*{D2: The exchange-rate identity and its concentration (T2, P2.1, P2.2)}

This subsection records the derivations behind T2 (the exchange-rate identity in
pathwise form), P2.1 (its concentration), and P2.2 (the feedback-bias constant).
Status of record: the identity was verified to zero relative error by a pre-existing
check; the pathwise refinement, the concentration rate, and the feedback-bias
constant are derived here and verified in V2. The formal statements live in the
register of theorem statements, and the proofs in the proofs appendix.

\paragraph{Setup.}
The stationary jittered loop is
\begin{equation*}
\dev_{t+1} \;=\; -m\,\dev_t + \sigma_e\,\xi_t,
\qquad
v \;=\; \Var(\dev_t) \;=\; \frac{\sigma_e^2}{1-m^2}.
\end{equation*}
OLS of the response $\tau_t$ on the quote $h_t$ has conditional-on-design variance
$\Var(\hat\eps \mid \mathrm{design}) = \sigma^2/\Sxx$ with
$\Sxx = \sum_{t\le T}(\dev_t-\bar{\dev})^2$, and the cost of Lemma L1 is
$\tfrac12\gpo\sum_{t\le T}\dev_t^2$.

\paragraph{The identity, upgraded from expectation to pathwise (T2).}
\textbf{Proposition (pathwise identity).} On every realization,
\begin{align*}
\Var(\hat\eps \mid \mathrm{design}) \times \tfrac12\gpo\sum_{t\le T}\dev_t^2
&\;=\; \tfrac12\,\gpo\,\sigma^2\,\frac{\sum_{t\le T}\dev_t^2}{\Sxx} \\
&\;=\; \tfrac12\,\gpo\,\sigma^2\Big(1+\frac{T\,\bar{\dev}^2}{\Sxx}\Big).
\end{align*}
Since $\bar{\dev}\to 0$ almost surely at rate $T^{-1/2}$ in the stationary regime,
the product equals $\tfrac12\gpo\sigma^2\,(1+O_P(1/T))$ pathwise: a strictly
stronger statement than the in-expectation identity, and the form the experiments
actually measure. The correction term $T\bar{\dev}^2/\Sxx$ is nonnegative, so the
realized product can only sit above the constant, by a vanishing amount.

\textbf{Invariance.} Neither $\sigma_e$, nor $m$, nor $T$ appears: more jitter buys
information and burns value at the same rate; a larger modulus amplifies the jitter
into more excitation and more cost at the same rate. The constant
$\tfrac12\gpo\sigma^2$ is the exchange rate between self-knowledge and value.

\paragraph{Concentration: the identity is observable (P2.1).}
The realized value cost (Lemma L1) carries the linear-term fluctuation
$-\delta_0\sum_{t\le T}\dev_t$. In the stationary regime $\sum_t\dev_t$ is a
mean-zero weakly dependent sum, so
\begin{equation*}
\frac{\operatorname{sd}(\text{realized } C_T)}{\E[C_T]} \;=\; O\big(T^{-1/2}\big),
\end{equation*}
with the explicit leading constant given by the AR(1) long-run variance
$v\,(1+\rho)/(1-\rho)$ at lag-one autocorrelation $\rho=-m$, that is
$v\,(1-m)/(1+m)$:
\begin{equation*}
\Var\Big(\sum_{t\le T}\dev_t\Big)
\;=\; T\,v\,\frac{1-m}{1+m}\,(1+o(1)).
\end{equation*}
The negative autocorrelation of the cobweb shrinks the fluctuation below the iid
level $T\,v$: a small bonus worth one sentence in the paper. Hence the measured
product concentrates around the invariant at rate $T^{-1/2}$. Verified in V2: the
seed-variance of the measured product scales as $1/T$ (it halves when $T$ doubles,
a factor-4 drop when $T$ quadruples, within Monte-Carlo tolerance).

\paragraph{Detail: $\Sxx$ versus $\sum_t\dev_t^2$.}
The estimator uses the centered $\Sxx$; the cost uses the raw $\sum_t\dev_t^2$.
Their ratio is
\begin{equation*}
\frac{\sum_{t\le T}\dev_t^2}{\Sxx}
\;=\; 1+\frac{T\,\bar{\dev}^2}{\Sxx}
\;=\; 1+O_P(1/T),
\end{equation*}
as above. Statements in the paper use whichever form is natural and record this
$O(1/T)$ equivalence once, in the present form.

\paragraph{The feedback (predetermined-regressor) bias, with its constant (P2.2).}
In the real loop the observed response feeds the next deployment. Model this as
\begin{equation*}
\dev_{t+1} \;=\; -m\,\dev_t + \sigma_e\,\xi_t + \phi\,\zeta_t,
\end{equation*}
with $\phi$ the noise-feedback gain: retraining on the realized flow moves the next
quote by $\phi\,\zeta_t$. Now $\dev_t$ is correlated with past $\zeta$, so the
regressor is predetermined but not strictly exogenous.

\textbf{Proposition (Stambaugh-type bias, full form, corrected by verification).}
The OLS slope satisfies
\begin{equation*}
\E\big[\hat\eps-(-\eps)\big]
\;=\; \frac{\phi\,\sigma^2\,(3m-1)}{(1-m^2)\,T\,v_\phi}\,(1+o(1)),
\qquad
v_\phi \;=\; \frac{\sigma_e^2+\phi^2\sigma^2}{1-m^2},
\end{equation*}
where $-\eps$ is the true slope in this sign convention (written
$\eps_{\mathrm{true}}$ in P2.2) and $v_\phi$ is the stationary variance of the
feedback loop. The bias is $O(1/T)$, proportional to the feedback gain, and changes
sign at $m=1/3$: slow-contracting loops bias the slope estimate upward,
fast-contracting ones downward, and at $m=1/3$ exactly the feedback bias vanishes
at first order.

\textbf{Derivation, both terms.} (Honesty record, kept per the falsification
convention: the first draft retained only the centering term and was falsified by
V2; the covariance term below is the correction.) Write $\hat b - b = N/S$ with
$N=\sum_{t\le T}(\dev_t-\bar{\dev})\,\zeta_t$ and $S=\Sxx$, and expand
\begin{equation*}
\E\Big[\frac{N}{S}\Big]
\;=\; \frac{\E[N]}{\E[S]} \;-\; \frac{\Cov(N,S)}{\E[S]^2} \;+\; O(T^{-2}).
\end{equation*}

\emph{Centering term.} Here $\E[\dev_t\zeta_t]=0$, since $\zeta_t$ is independent
of the past-built $\dev_t$, so $\E[N]=-\E\big[\bar{\dev}\sum_t\zeta_t\big]$. With
$\E[\dev_s\zeta_t]=\phi\sigma^2(-m)^{s-t-1}$ for $s>t$,
\begin{equation*}
\E[N]
\;=\; -\frac{1}{T}\sum_{s>t}\phi\,\sigma^2\,(-m)^{s-t-1}
\;=\; -\frac{\phi\,\sigma^2}{1+m} \;+\; O(1/T).
\end{equation*}

\emph{Covariance term (the one the first draft missed).} By the Gaussian
fourth-moment (Wick) expansion,
\begin{equation*}
\Cov(N,S)
\;=\; 2\sum_{t,s}\E[\tilde{\dev}_t\tilde{\dev}_s]\,\E[\tilde{\dev}_s\zeta_t],
\qquad \tilde{\dev} = \dev-\bar{\dev};
\end{equation*}
substituting $\E[\tilde{\dev}_t\tilde{\dev}_s]\simeq v_\phi\,(-m)^{|s-t|}$ and the
cross-covariance above,
\begin{equation*}
\Cov(N,S)
\;=\; 2\,v_\phi\,\phi\,\sigma^2\,T\sum_{k\ge 1}(-m)^{k}(-m)^{k-1}
\;=\; -\,\frac{2\,v_\phi\,\phi\,\sigma^2\,T\,m}{1-m^2}.
\end{equation*}

\emph{Combine.} With $\E[S]\simeq T\,v_\phi$,
\begin{equation*}
\mathrm{bias}
\;=\; \Big[-\frac{1}{1+m}+\frac{2m}{1-m^2}\Big]\,
\frac{\phi\,\sigma^2}{T\,v_\phi}
\;=\; \frac{3m-1}{1-m^2}\cdot\frac{\phi\,\sigma^2}{T\,v_\phi},
\end{equation*}
which is P2.2. Verified in V2 at two moduli straddling the sign change: $m=0.5$
gives a positive bias and $m=0.2$ a negative one, each within Monte-Carlo tolerance
of the constant. The sign change is the sharpest possible falsification test of the
two-term structure, and it passes.

\textbf{Paper usage.} State the identity (T2) under exploration-dominant design;
report the bias formula (P2.2) as the known, computable deviation, a checkable
condition ($\phi$ is estimable from the loop's own update rule), not an assumption.

\paragraph{Verified numerically (V2).}
The V2 checks for this subsection: (1) the pathwise product equals
$\tfrac12\gpo\sigma^2\,(1+T\bar{\dev}^2/\Sxx)$ exactly; (2) product concentration:
the seed-variance of the measured product scales as $1/T$, a factor-4 drop when $T$
quadruples, within Monte-Carlo tolerance; (3) the long-run variance of
$\sum_t\dev_t$ matches $T\,v\,(1-m)/(1+m)$; (4) the feedback bias: sign, magnitude,
and $1/T$ scaling of $\phi\,\sigma^2\,(3m-1)/\big((1-m^2)\,T\,v_\phi\big)$.

\subsection*{D3: Minimax lower bounds and the exact reach (T3, L2, T4, T9)}

This record derives the two lower bounds that turn the exchange-rate identity
of D2 (register T2) from a property of one scheme into a property of the
problem, and then the exact reach within which that floor is structure-proof
against a policy that knows the response family (T9). T2 shows that a
specific scheme, stationary symmetric jitter with ordinary least squares,
achieves $\Var(\hat\eps)\times\text{cost} = \tfrac{1}{2}\gpo\sigma^2$. The
claim that makes this an exchange rate is that no adaptive exploration
policy and no estimator does better. There are two versions of the bound,
one per budget type: an exact deviation-budget bound (T3, derived in full
and verified numerically in V3a) and a value-budget bound (T4), whose engine
is the exploitation--information lemma (L2), proved below with its
$O(T^{-1/2})$ correction established and verified in V3b. Throughout, the
formal proofs are recorded in the proofs appendix; this subsection records
the derivations and the caveats attached to them.

\paragraph{Deviation budget: van Trees over adaptive designs (T3).}
Setting: the exploration policy is arbitrary and adaptive (each $u_t$ any
measurable function of the history), subject to the pathwise budget
$\sum_{t\le T}\dev_t^2 \le D$ almost surely; the estimator $\hat\eps$ of
$\eps$ is arbitrary, bias allowed. The claim (T3) is that the minimax risk
over an interval of $\eps$ values of width $w$ satisfies
\begin{equation*}
\inf_{\text{policy, estimator}}\;\sup_{\eps}\;
\E\big[(\hat\eps-\eps)^2\big]
\;\ge\; \frac{\sigma^2}{D+\sigma^2(\pi/w)^2}
\;=\; \frac{\sigma^2}{D}\,\big(1-o(1)\big)
\qquad\text{as } D\,(w/\sigma)^2\to\infty .
\end{equation*}
Multiplying by the budget cost $\tfrac{1}{2}\gpo D$ (cost equivalence, L1)
gives
\begin{equation*}
\sup_{\eps}\;\Var(\hat\eps)\times\tfrac{1}{2}\gpo D
\;\ge\; \tfrac{1}{2}\,\gpo\,\sigma^2\,\big(1-o(1)\big):
\end{equation*}
the exchange rate is a floor.

The derivation is the van Trees inequality (Bayesian Cram\'er--Rao, in the
form of Gill and Levit, 1995) with a prior density $\pi_0$ supported on the
width-$w$ interval:
\begin{equation*}
\E_{\pi_0}\,\E\big[(\hat\eps-\eps)^2\big]
\;\ge\; \frac{1}{\E_{\pi_0}\big[I_T(\eps)\big] + I(\pi_0)},
\qquad
I(\pi_0)=\int \frac{(\pi_0')^2}{\pi_0},
\end{equation*}
where $I(\pi_0)$ is the prior information, equal to $(\pi/w)^2$ for the
cosine-squared prior (the minimizer of prior information on an interval of
width $w$), and $I_T(\eps)$ is the Fisher information of the adaptively
designed experiment. The one nonstandard step is that $I_T$ for an adaptive
design is still bounded by the design energy: by the chain rule for Fisher
information over the filtration,
\begin{equation*}
I_T(\eps)
\;=\; \sum_{t\le T}\E\big[I(\tau_t \mid \mathcal{F}_t;\eps)\big]
\;=\; \sum_{t\le T}\frac{\E[\dev_t^2]}{\sigma^2},
\end{equation*}
because conditionally on $\mathcal{F}_t$ the observation $\tau_t$ is Gaussian
with mean linear in $\eps$ (slope $-\dev_t$, with $\dev_t$
$\mathcal{F}_t$-measurable): the policy can place information but cannot
manufacture it. The pathwise budget then gives $\E[I_T]\le D/\sigma^2$;
substituting, and taking the supremum over the interval
($\sup \ge \E_{\pi_0}$), yields the bound.

Three remarks recorded with the derivation. (i) Biased estimators are
covered: van Trees needs no unbiasedness. (ii) The bound is achieved, up to
the $o(1)$, by two-point symmetric probing at the budget with least squares,
tying back to the D2 identity (T2); this is the ``attained with equality''
clause of T3. (iii) The $o(1)$ is explicit: it equals
$\sigma^2(\pi/w)^2/D$.

\paragraph{D3b, the adaptive companion. Value budget: the exploitation problem.}
Under the value budget the policy pays
$\Phi(\hstar_{\mathrm{CE}})-\Phi(h_t)$ per step (the certainty-equivalent
anchor of D0). The residual linear coefficient at the anchor is
\begin{equation*}
\varphi'(\bar\eps)\,(\eps-\bar\eps)
\;=\; -\beta(\hstar-\psi)(\eps-\bar\eps),
\end{equation*}
so a policy could try to earn while exploring: guess the sign of
$\eps-\bar\eps$ and drift accordingly. The lemma below says this
self-financing is bounded by the information already purchased: the policy
cannot exploit what it has not yet paid to know.

Setting: symmetric two-point prior
$\eps\in\{\bar\eps-\delta,\ \bar\eps+\delta\}$; write
$s=\operatorname{sign}(\eps-\bar\eps)$ and
$g_1 := \beta\,|\hstar-\psi|$, the known coefficient of the unknown part.
The policy is non-anticipating, with per-step trust region $|\dev_t|\le r$;
write $S_t := \sum_{i\le t}\dev_i^2$ for the design energy.

\paragraph{The exploitation--information lemma (L2).}
Two ingredients. First, the posterior-imbalance step: the policy's knowledge
of $s$ at time $t$ obeys
\begin{equation*}
\big|\E[s\mid\mathcal{F}_t]\big| \;\le\; \TV(P^+_t,P^-_t) \;=:\; \TV_t
\end{equation*}
(posterior imbalance is bounded by the distinguishability of the two
environments; the starting prior is symmetric). Register L2 records the
averaged form of this step as the exact identity
$\E\big|\E[s\mid\mathcal{F}_t]\big| = \TV(P^+_t,P^-_t)$. Second, the
Pinsker step: the two environments differ only in the observation means, by
$2\delta\dev_i$ at step $i$, so
\begin{equation*}
\KL_t \;=\; \sum_{i\le t}\frac{(2\delta\dev_i)^2}{2\sigma^2}
\;=\; \frac{2\delta^2 S_t}{\sigma^2},
\qquad
\TV_t \;\le\; \sqrt{\KL_t/2} \;=\; \frac{\delta\sqrt{S_t}}{\sigma}.
\end{equation*}
(For adaptive designs the same computation runs through the chain rule for
$\KL$ over the filtration, with $S_t$ replaced by $\E S_t$, exactly as in the
Fisher computation for T3; this is the form registered in L2.)

Chaining these, the expected exploitation gain obeys, for any
non-anticipating policy,
\begin{align*}
\E[\mathrm{gain}_T]
\;:=\; \E\Big[\sum_{t\le T} g_1\,\delta\, s\,\dev_t\Big]
&\;\le\; g_1\,\delta\,\sum_{t\le T}\E\big[\TV_t\,|\dev_t|\big] \\
&\;\le\; \frac{g_1\,\delta^2}{\sigma}\,\sqrt{S_T}\,
\sum_{t\le T}\E\,|\dev_t|
\;\le\; \frac{g_1\,\delta^2}{\sigma}\, S_T\,\sqrt{T},
\end{align*}
where the first inequality is the tower rule plus the posterior-imbalance
step, the second uses $\TV_t\le\delta\sqrt{S_t}/\sigma
\le\delta\sqrt{S_T}/\sigma$, and the last is Cauchy--Schwarz
($\sum_{t\le T}|\dev_t| \le \sqrt{T\,S_T}$). For fully adaptive designs the
same chain, run with expected energies, gives the closed form recorded in
L2, $\E[\mathrm{gain}_T]\le g_1\,\delta^{3/2}\,T^{1/2}
(\E S_T)^{3/4}/\sigma^{1/2}$; both forms are $O(\sigma\sqrt{T})$ at the
minimax scale below. The formal proof is recorded in the proofs appendix
under L2.

\paragraph{The minimax scale and the value-budget bound (T4).}
The minimax-relevant prior scale is
$\delta = c\,\sigma/\sqrt{S_T}$: this is the accuracy the budget can buy
(any larger $\delta$ is learned exactly, any smaller is unlearnable).
Substituting into the lemma,
\begin{equation*}
\E[\mathrm{gain}_T] \;\le\; g_1\, c^2\,\sigma\,\sqrt{T},
\qquad\text{while}\qquad
\E[C_T] \;=\; \tfrac{1}{2}\,\gpo\, S_T \;\sim\; \tfrac{1}{2}\,\gpo\, T\, v
\end{equation*}
(cost equivalence L1, stationary-scale exploration $S_T=\Theta(T)$ with
per-step exploration variance $v$). The exploitable fraction of the budget
is therefore
\begin{equation*}
\frac{\E[\mathrm{gain}_T]}{\E[C_T]}
\;\le\; \frac{2\,g_1\,c^2\,\sigma}{\gpo\, v\,\sqrt{T}}
\;=\; O\big(T^{-1/2}\big).
\end{equation*}
The Le Cam two-point reduction then quantifies the constant. At the minimax
scale $\TV_T\le c$, so any estimator's worst-case risk over the pair obeys
\begin{equation*}
\sup_{\eps}\,\E\big[(\hat\eps-\eps)^2\big]
\;\ge\; \frac{\delta^2}{2}\,\big(1-\TV_T\big)
\;\ge\; \frac{\delta^2}{2}\,(1-c).
\end{equation*}
Multiplying by $\E[C_T]=\tfrac{1}{2}\gpo S_T$, using
$\delta^2 S_T = c^2\sigma^2$, and applying the exploitation rebate (the
value budget can be underpaid only by the exploited gain, an
$O(T^{-1/2})$ fraction by the display above):
\begin{equation*}
\sup\;\Var(\hat\eps)\times C_T
\;\ge\; \frac{\gpo\,\sigma^2}{4}\; c^2\,(1-c)\;
\big(1-O(T^{-1/2})\big).
\end{equation*}
The factor $c^2(1-c)$ is maximized at $c=2/3$, with value $4/27$, giving the
registered bound of T4:
\begin{equation*}
\inf_{\text{policy, estimator}}\;\sup\;\Var(\hat\eps)\times C_T
\;\ge\; \frac{\gpo\,\sigma^2}{27}\,\big(1-O(T^{-1/2})\big).
\end{equation*}
The sharp constant $\tfrac{1}{2}$, matching the achievable product of T2,
that is
$\inf\sup\,\Var(\hat\eps)\times C_T
\ge \tfrac{1}{2}\gpo\sigma^2\,(1-O(T^{-1/2}))$,
is conjectured and is part of OPEN-1: no simulated policy fell below it by
more than Monte-Carlo error (\texttt{run\_open1.py}, part B).

\paragraph{Interpretation: the self-referential structure.}
The profitable direction is $s$, the very thing being estimated. Pinsker
converts ``how much the policy knows about $s$'' into ``how much design
energy it has already spent''; any gain is therefore a rebate proportional
to information already paid for, and the rebate's rate vanishes. Knowledge
cannot be bootstrapped from its own purchase. To our knowledge no analog of
this lemma appears in the performative-prediction, pricing, or
experimental-design literatures (literature review items G1 to G3).

\paragraph{Caveats recorded with the derivation.}
(i) The constant in the $O(T^{-1/2})$ term is explicit,
$2 g_1 c^2 \sigma/(\gpo v)$, and blows up as $v\to 0$: a policy that
barely explores can have its tiny cost substantially rebated. This is
consistent, since in that regime both sides of the product are dominated by
the transient (D1); the theorem is stated for stationary-scale exploration,
$S_T=\Theta(T)$.
(ii) $g_1=\beta|\hstar-\psi|$ is $O(1)$ in general; in the
weak-performativity regime it is $O(\eps)$, giving the sharper factor
$\big(1-O(\eps\, T^{-1/2})\big)$.
(iii) The fixed-$\delta$ regime is different, and the certainty-equivalent
anchor is the reason. At fixed prior separation the policy eventually learns
$s$ outright and can exploit it indefinitely: the exploitation fraction does
not vanish (observed directly in V3b's first, deliberately mis-scaled run).
That regime is not a counterexample. Sustained exploitation of learned
structure is exactly what D0's certainty-equivalent anchor re-anchors away:
once $\eps$ is known, the CE-optimal deployment shifts, and the ``gain'' is
no longer measured as a rebate against exploration. The theorem's content
concerns the unlearned residual, whose hardest instances live at the minimax
scale $\delta\sim\sigma/\sqrt{S_T}$, and there the rebate vanishes at rate
$T^{-1/2}$ (verified). We state this distinction explicitly: it separates
the theorem's content from an apparent counterexample.

\paragraph{Numerical verification (V3).}
Three checks. (V3a) A Bayes-risk simulation with a Gaussian-prior
posterior-mean estimator under three designs (front-loaded, spread,
adaptive-greedy) with capped budget $D$: the measured risk respects
$\sigma^2/\big(D+\sigma^2/\sigma_{\pi}^2\big)$ in every arm, with the spread
design near-tight. (V3b-i) The Pinsker chain: the measured posterior
imbalance satisfies
$\big|\E[s\mid\mathcal{F}_t]\big| \le \delta\sqrt{S_t}/\sigma$ pathwise
along simulated runs. (V3b-ii) An explicitly exploiting policy, which
drifts by the posterior mean of the unknown linear term, simulated at
several horizons $T$: its realized $\mathrm{gain}_T/C_T$ decays
consistently with $O(T^{-1/2})$, and its achieved
$\Var(\hat\eps)\times\text{cost}$ never falls below
$\tfrac{1}{2}\gpo\sigma^2\,(1-\text{measured rebate})$.

\paragraph{Structure-proofness within reach (T9).}
The floors above are nonparametric. A separate threat is a policy that
knows the response family's functional form $\tau(h)=C_0+C_1 e^{-ch}$ and
shops for designs whose parametric Cram\'er--Rao product undercuts the
floor. Define, for a finitely supported design $\mu$,
\begin{equation*}
R(\mu) \;=\; \big(g^\top M(\mu)^{-1} g\big)\!\int (h-\hstar)^2\, d\mu,
\qquad
M(\mu) \;=\; \int s\, s^\top d\mu,
\end{equation*}
with the sensitivity and the estimand gradient
\begin{equation*}
s(h) \;=\; \big(1,\; e^{-ch},\; -C_1\, h\, e^{-ch}\big)^\top,
\qquad
g \;=\; \nabla_\theta\, \eps(\hstar) \;=\; -s'(\hstar),
\end{equation*}
normalized so that $R(\mu)\ge 1$ says the family-knowing product is at
least the exchange-rate floor $\tfrac{1}{2}\gpo\sigma^2$ under the
local-quadratic (A1) cost.

\paragraph{One-dimensional reduction.}
The generalized Rayleigh identity plus the bijection
$u \leftrightarrow \phi_u = u^\top s$ between $\mathbb{R}^3$ and
$V=\operatorname{span}\{1,\ e^{-ch},\ h e^{-ch}\}$ (using $C_1\neq 0$)
gives
\begin{equation*}
g^\top M(\mu)^{-1} g
\;=\; \frac{1}{\inf\big\{\ \|\phi\|_\mu^2 \ :\ \phi\in V,\
\phi'(\hstar)=1\ \big\}},
\end{equation*}
because $u^\top g = -\phi_u'(\hstar)$ (the sign squares away). So the floor
holds at $\mu$ if and only if some unit-slope family element is
$\mu$-cheaper than the linear function $\dev(h)=h-\hstar$.

\paragraph{The candidate and its reach.}
The unique element of $V$ with 2-jet $(0,1,0)$ at the anchor is
\begin{equation*}
\phi_0(h) \;=\; \frac{2}{c}
\;-\; e^{-c(h-\hstar)}\Big(\frac{2}{c} + (h-\hstar)\Big),
\end{equation*}
with expansion $\phi_0(x) = x - (c^2/6)\,x^3 + \cdots$ in $x=h-\hstar$.
The pointwise domination $|\phi_0(x)|\le |x|$ holds exactly on
$x\in[-2/c,\ \infty)$ (equality only at $x=0$ and $x=-2/c$), and fails
strictly and exponentially below $-2/c$. The three elementary ingredients
are $1+t < e^t$ (which makes $\phi_0(x)-x$ strictly decreasing),
$\phi_0'(x) = e^{-cx}(1+cx) > 0$ above $-1/c$, and the single sign change
of $e^{-cx}(1+cx)+1$ on $(-2/c,\,0)$ (since $e^{-cx}(1+cx)$ is increasing
there from $-e^2$ to $1$). The full proof is recorded in the proofs
appendix under T9; verified in V9.1.

\paragraph{Consequence (T9(i)).}
For any design supported in $[\hstar-2/c,\ \infty)$: $R(\mu)\ge 1$
whenever the estimand is estimable, strict as soon as $\mu$ places mass
off the two equality points $\{\hstar,\ \hstar-2/c\}$, and
$\inf_\mu R(\mu)=1$, approached by symmetric
three-point designs collapsing at the anchor (at rate $\delta^2$ in the
collapse half-width $\delta$, V9.2; the limit is exact because in 2-jet
coordinates the collapsing model is the quadratic family, where every
symmetric design has $R=1$ on the nose). The trust region A4 with
$r\le 2/c$ is therefore exactly the condition under which the floor is
structure-proof, not a technicality.


\paragraph{Sharpness of the candidate, and an exact duality.}
Three facts sharpen the picture, established in the adversarial
verification campaign with certified numerics (the full write-up is
journal material under OPEN-5). (i) The easy direction of a duality: for
a finite support $X$ let
$m(X) = \max\{\phi'(\hstar) : \phi \in V,\ |\phi(x)| \le |x - \hstar|
\ \forall x \in X\}$; then for \emph{every} weighting $w$ of $X$,
$R(X, w) \ge m(X)^2$, because the maximizer $\psi^*$ scaled to unit slope
is feasible in the Rayleigh infimum with
$\int (\psi^*/m)^2\, dw \le \int d^2\, dw / m(X)^2$. Within reach
$\phi_0$ certifies $m(X) \ge 1$, which is exactly T9(i); numerically the
duality is tight ($\inf_w R(X, w) = m(X)^2$ by linear-programming
duality, certified on random supports), so pointwise domination is not a
lossy proof device: the reach question \emph{is} the domination question.
(ii) Single-candidate sharpness: among all unit-slope elements of $V$,
$\phi_0$ is the unique one whose domination set contains a full interval
$[\hstar - a, \hstar)$ with $a > 0$, and the maximal such $a$ is exactly
$2/c$. The affine family with 2-jet $(0, 1, \gamma)$, $\gamma \in (-1, 0)$,
trades the reach boundary outward (to $u_2(\gamma) < -2/c$ in units of
$1/c$) at the proved price of vacating an interval just below the anchor,
and the violating designs live exactly in the vacated gap. (iii) The
below-reach failure is generic, not a boundary artifact: for every
$A > 2/c$ there are nonsingular designs supported in
$[\hstar - A, \infty) \cap (0, \infty)$ with $R < 1$ (a certified
three-point example in the register market attains $R = 0.8537$).
Consequence for the statement of T9: by the extended-Chebyshev property
of $V$ (any nonzero element has at most two zeros counting multiplicity),
every design with at least three distinct support points automatically
has nonsingular information, so the estimability caveat in T9 concerns
only one- and two-point designs.

\paragraph{The reach is exact (T9(ii)).}
Below $\hstar-2/c$ the domination fails, and the failure is realized: the
frozen witness (support $\{1.8726,\allowbreak\ 1.8665,\allowbreak\ 1.3073,\allowbreak\ 0.05\}$, weights
$\{0.003329,\allowbreak\ 0.955702,\allowbreak\ 0.040697,\allowbreak\ 0.000272\}$ normalized, frozen anchor
$\hstar=1.8461$, register market $c=1.5$) attains
\begin{equation*}
R \;=\; 0.8516504721831888870\ldots
\end{equation*}
(50-digit precision; V9.2), with every support point except the $0.05$
probe inside the reach: one vanishing-weight below-reach probe breaks the
floor by 15 percent. Its true-cost ratio is $1.011>1$, so the violation
lives in the local-quadratic cost model, which is the floor's own scope.
Symmetric pair-plus-far-probe families never violate (their ratio tends to
$1$ from above as the probe weight vanishes; the old scan's verdict of
global structure-proofness was a weight-grid artifact, recorded as the
eleventh measurement-forced pivot). The characterization of the full
below-reach violation set is OPEN-5.

\subsection*{D4/D5: Design geometry: optimal shapes, dispersion, temporal shaping, Chebyshev support (T5, C5.1, L3, T6)}

This subsection records the derivations behind the three exact design optima
of T5, the dispersion corollary C5.1, the temporal-shaping collapse L3, and
the Chebyshev counting result T6; it covers both register slots D4 and D5,
the Chebyshev material being the D5 portion. All three design optima are
derived with proofs, the temporal-shaping lemma and the counting result are
proved, and everything is verified in the V4 and V5 numerical checks. It is
a derivation record: each optimum is computed in full, the caveats are
stated where they arose, and the numerical anchors are those of V4 and V5.
The formal statements live in the theorem register under the IDs above.

\paragraph{Setting.}
The decision is $d$-dimensional and the estimand is the response Jacobian.
Exploration has stationary second moment $M := \E[\dev\dev^\top]$ (positive
semidefinite), the response noise has variance $\sigma^2$ per coordinate,
and by the $d$-dimensional form of the cost equivalence lemma L1 (from D0)
the value cost rate is
\begin{equation*}
\tfrac{1}{2}\,\E\big[\dev^\top \Gpo\, \dev\big]
\;=\; \tfrac{1}{2}\,\tr(\Gpo M),
\end{equation*}
so the value budget over horizon $T$ is $B = (T/2)\,\tr(\Gpo M)$. The
per-unit-time information matrix is $M/\sigma^2$. Absorbing $T$ and
$\sigma^2$ into units, all three problems take the static normal form of a
linear budget constraint $\tr(\Gpo M) \le B$ on the design $M$.

\paragraph{A-optimality (T5(a)): what to do if you must know everything
equally.}
The problem and its exact solution are
\begin{align*}
& \text{minimize} \quad \tr\big(M^{-1}\big)
\qquad \text{subject to} \quad \tr(\Gpo M) \le B, \quad M \succ 0, \\
& M^{*} \;=\; \frac{B}{\tr\,\Gpo^{1/2}}\;\Gpo^{-1/2},
\qquad
\tr\big((M^{*})^{-1}\big) \;=\; \frac{\big(\tr\,\Gpo^{1/2}\big)^{2}}{B}.
\end{align*}
Derivation: the objective is strictly convex on the positive definite cone
and the constraint is linear, so the KKT conditions are sufficient.
Stationarity reads $-M^{-2} + \lambda\,\Gpo = 0$, hence
$M = \lambda^{-1/2}\,\Gpo^{-1/2}$; the budget pins the multiplier through
$\tr(\Gpo M) = \lambda^{-1/2}\,\tr\,\Gpo^{1/2} = B$, so
$\lambda^{-1/2} = B/\tr\,\Gpo^{1/2}$, and substituting back gives the
stated value. A diagonal-case Cauchy--Schwarz check confirms the shape:
with curvatures $g_i$ and design weights $m_i$,
\begin{equation*}
\Big(\sum_i \frac{1}{m_i}\Big)\Big(\sum_i g_i\, m_i\Big)
\;\ge\; \Big(\sum_i \sqrt{g_i}\Big)^{2},
\end{equation*}
with equality at $m_i \propto g_i^{-1/2}$.

The reading is: explore where the objective is flat, but only to the
square-root degree. The optimal exploration covariance is $\Gpo^{-1/2}$,
the inverse square root and not the inverse: flat directions are cheap per
unit information but not infinitely favoured, while curved directions are
expensive but cannot be abandoned when every coordinate of the Jacobian is
needed.

\paragraph{The price of isotropic jitter (C5.1).}
Isotropic exploration at the same budget, $M = (B/\tr\,\Gpo)\,I$, is
feasible with the budget tight and has risk
$\tr(M^{-1}) = d\,\tr(\Gpo)/B$. The ratio to the A-optimal value is
\begin{equation*}
F \;=\; \frac{d\,\tr(\Gpo)}{\big(\tr\,\Gpo^{1/2}\big)^{2}} \;\ge\; 1,
\end{equation*}
with equality if and only if all curvatures are equal (Cauchy--Schwarz
again). $F$ is the curvature dispersion of the objective: the exact factor
by which naive isotropic exploration, which is optimal in LQR-like settings
(Simchowitz and Foster, 2020), overpays in the performative setting. This
is the content of C5.1, restated as the isotropy contrast P9.2, and $F$ is
computable on the calibrated multi-bond $\Gpo$ (whose diagonal is given
exactly by P9.3).

\paragraph{D-optimality (T5(b)): what the confidence-volume objective
wants.}
\begin{align*}
& \text{maximize} \quad \log\det M
\qquad \text{subject to} \quad \tr(\Gpo M) \le B, \\
& M^{*} \;=\; \frac{B}{d}\;\Gpo^{-1},
\qquad
\log\det M^{*} \;=\; d\,\log\frac{B}{d} \;-\; \log\det \Gpo .
\end{align*}
Derivation: $\nabla \log\det M = M^{-1}$, so stationarity reads
$M^{-1} = \lambda\,\Gpo$; the objective is concave and the constraint
linear, and the budget gives $\lambda^{-1} = B/d$.

A caveat recorded at derivation time: the two shapes differ, A-optimality
giving $\Gpo^{-1/2}$ and D-optimality giving $\Gpo^{-1}$, and the two must
never be conflated in the paper. The algorithm derivation (D6) uses the
D-optimal criterion internally (Frank--Wolfe on $\log\det$), while the
headline risk statement is A-optimal; both are stated, and the discrepancy
factor between them is itself computable and is small when the dispersion
$F$ is moderate.

\paragraph{c-optimality (T5(c)): what the corrected (PerfGD-style) update
actually needs.}
This is the new result of the three. The corrected update consumes the
response estimate through a single functional, the correction direction
$c$: in the scalar family the correction $\Delta$ needs only
$\eps(\hstar)$, and in $d$ dimensions $c = \nabla_{E}\,\Delta$ is a fixed
vector once the operating point is known. Estimating the scalar
$c^\top\theta$ alone:
\begin{align*}
& \text{minimize} \quad c^\top M^{-1} c
\qquad \text{subject to} \quad \tr(\Gpo M) \le B, \\
& M^{*} \;=\; B\,\frac{c\,c^\top}{c^\top \Gpo\, c}
\quad \text{(rank one: probe along $c$ itself)},
\qquad
c^\top (M^{*})^{+}\, c \;=\; \frac{c^\top \Gpo\, c}{B}.
\end{align*}
Derivation: substitute $N = \Gpo^{1/2} M\, \Gpo^{1/2}$, so that
$\tr N = \tr(\Gpo M) \le B$, and $w = \Gpo^{1/2} c$; the objective becomes
$w^\top N^{-1} w$. By Cauchy--Schwarz,
\begin{equation*}
\big(w^\top w\big)^{2}
\;=\; \big(w^\top N^{-1/2}\, N^{1/2}\, w\big)^{2}
\;\le\; \big(w^\top N^{-1} w\big)\big(w^\top N w\big),
\end{equation*}
and $w^\top N w \le \|w\|^{2}\,\tr N \le \|w\|^{2} B$, so
$w^\top N^{-1} w \ge \|w\|^{2}/B$, attained at
$N^{*} = B\,ww^\top/\|w\|^{2}$. Back-substituting,
\begin{equation*}
M^{*} \;=\; \Gpo^{-1/2} N^{*}\, \Gpo^{-1/2}
\;=\; B\,\frac{\big(\Gpo^{-1/2} w\big)\big(\Gpo^{-1/2} w\big)^\top}
{\|w\|^{2}},
\qquad \Gpo^{-1/2} w = c,\quad \|w\|^{2} = c^\top\Gpo\,c,
\end{equation*}
giving $M^{*} = B\,cc^\top/(c^\top \Gpo\, c)$ as stated.

The corrected trio is worth displaying side by side:
\begin{equation*}
\text{A-opt:}\;\; M^{*} \propto \Gpo^{-1/2},
\qquad
\text{D-opt:}\;\; M^{*} \propto \Gpo^{-1},
\qquad
\text{c-opt:}\;\; M^{*} \propto c\,c^\top .
\end{equation*}
A-optimality shapes exploration as $\Gpo^{-1/2}$, D-optimality as
$\Gpo^{-1}$, and c-optimality probes along the correction direction itself:
the curvature does not tilt the direction at all, it only prices it. The
value $c^\top \Gpo\, c / B$ says that the cost of knowing your correction
is the $\Gpo$-norm of what the correction asks for. The savings over the
A-optimal value is the operational argument for correction-targeted
probing: a desk does not need the whole Jacobian, and the geometry says
exactly what it may skip.

Two caveats, recorded as they arose. First, an earlier draft of this
derivation mis-stated the c-optimal direction as $\Gpo^{-1} c$; the
substitution above corrects it. The error is recorded here per the
project's falsification convention, and the corrected direction is locked
by the V4 numerical check. Second, rank-one designs are singular (hence the
pseudo-inverse in the value); in practice they are regularized by the trust
region, and the paper states both the limiting and the regularized forms.

\paragraph{Temporal shaping is irrelevant (L3): design space collapse.}
\textbf{Lemma L3 (temporal shaping is irrelevant).} With response noise
i.i.d.\ and independent of the deployment path, the information about the
response parameters depends on the exploration process only through its
empirical second moment $\sum_t \dev_t \dev_t^\top$. The computation is one
line: the conditional log-likelihood is
\begin{equation*}
\sum_t -\frac{\big(\tau_t - a - \theta^\top \dev_t\big)^{2}}{2\sigma^{2}},
\end{equation*}
and its Hessian in $\theta$ is $-\sum_t \dev_t \dev_t^\top / \sigma^{2}$,
regardless of the temporal order or dependence structure of $\{\dev_t\}$.
Hence temporally correlated exploration (AR-shaped, cyclic, burst) buys
nothing beyond its stationary design measure, and the design problem is
exactly the static problems treated above. Scope boundary: the lemma fails
when the noise is serially correlated or feedback-coupled (the regime of
section 4 of D2), where shaping does matter; this is stated as the
boundary of L3 and left to the journal version.

\paragraph{Chebyshev counting (T6): why curvature is invisible to
trajectories.}
This is the D5 portion of the record. The structural family of REFLEX
theory 1.1 (\texttt{structural\_response.py}) is
\begin{equation*}
\tau(h) \;=\; C_0 + C_1\, e^{-c h},
\qquad \text{parameters } (C_0, C_1, c), \quad p = 3,
\end{equation*}
with sensitivity vector
\begin{equation*}
s(h) \;=\; \big(\,1,\; e^{-c h},\; -C_1\, h\, e^{-c h}\,\big)^{\!\top}.
\end{equation*}
(In this paragraph $c$ is the decay rate of the structural family, not the
correction direction of the c-optimal problem.) The counting argument has
four parts.

(i) $\{1,\, e^{-ch},\, h e^{-ch}\}$ is an extended Chebyshev system on any
interval: the relevant Wronskian-type determinants are nonzero (standard).

(ii) Consequently the Fisher matrix
$M(\xi) = \sum_i w_i\, s(h_i)\, s(h_i)^\top$ of any design $\xi$ with fewer
than three distinct support points is singular, and the decay rate $c$ is
unidentified. The computation: $\operatorname{rank} M(\xi) \le
\#\mathrm{support}(\xi)$, so a two-point design gives rank at most
$2 < 3 = p$, hence $\det M(\xi) = 0$ and the $c$-direction lies in (or
meets) the null space. The exact-arithmetic check in V5 finds
$\det M = 0$ to machine precision for every random two-point design and
$\det M > 0$ for generic three-point designs.

(iii) Any noiseless RRM trajectory has effective support at most two
(C1.2 of D1: geometric collapse for modulus below one, a period-two orbit
at the boundary). Hence no retraining trajectory, at any modulus,
identifies the curvature of its own response family.

(iv) D-optimal designs for this family on a trust-region interval
$[\hstar - r,\, \hstar + r]$ are supported on exactly three points
including both endpoints (a de la Garza-type argument; the interior point
is computed numerically in the experiments).

This is the sharp form of the repository's empirical rule that exponential
response fits need a wide spread range, and the formal reason the
anti-echo freeze (the \texttt{identified} flag in
\texttt{structural\_response.py}) exists.

\paragraph{Verified numerically (V4, V5).}
Six checks anchor the results above. (1)~A-optimality: random search over
positive semidefinite $M$ at fixed budget never beats
$(\tr\,\Gpo^{1/2})^{2}/B$, and the closed-form $M^{*}$ achieves it (run at
$d = 5$ with random symmetric positive definite $\Gpo$). (2)~D-optimality:
$M^{*} = (B/d)\,\Gpo^{-1}$ beats random feasible designs in $\log\det$.
(3)~c-optimality: random search never beats $c^\top \Gpo\, c / B$, and the
rank-one construction approaches it. (4)~The isotropic factor
$F = d\,\tr(\Gpo)/(\tr\,\Gpo^{1/2})^{2}$ matches the measured ratio.
(5)~Temporal lemma: AR-shaped and i.i.d.\ exploration with matched
empirical second moment give identical OLS covariance. (6)~Chebyshev:
$\det M = 0$ for two-point designs versus $\det M > 0$ for three-point
designs, and the $c$-direction variance is infinite or huge under
two-point designs.

\subsection*{D6: Safe certainty equivalence (L4, R2, P7.1)}

This derivation records the safety analysis of the certainty-equivalent
corrected loop: the structural observation R2 (open-loop exploration cannot
destabilize the linearized retraining dynamics), the perturbed-modulus
lemma L4 (how $C^1$ estimation error in the response moves the closed-loop
modulus), and the safety proposition P7.1 (the pessimism gate inherits
uniform-in-time safety from an anytime-valid confidence sequence). The
formal statements are recorded in the register under these IDs, and the
proofs in the proofs appendix. Status of the record: R2 and L4 are derived
and verified (V6); P7.1 is proved conditional on the validity of the
confidence sequence; the $O(\sqrt{T})$ design-regret bound in the closing
paragraph is stated with a proof strategy only and is labeled a
journal-version item.

\paragraph{Where the stability constraint actually bites (R2).}
A fact discovered in the derivation and promoted to the paper: open-loop
exploration cannot destabilize the linearized loop. The linearized
deviation dynamics under a planned (history-independent) probe sequence
$(u_t)$ read
\begin{equation*}
\dev_{t+1} \;=\; -m\,\dev_t \;+\; u_t ,
\end{equation*}
and the homogeneous part keeps its modulus $m<1$ regardless of the input
$u$: bounded inputs give bounded outputs. This is R2. The stability risk
of probing therefore enters through exactly two channels.

First, estimate feedback, the real one. The corrected update deploys
\begin{equation*}
h_{t+1} \;=\; h_t + \eta\,\hat\Phi'(h_t),
\qquad
\hat\Phi' \;=\; G + \hat\Delta ,
\end{equation*}
where $\hat\Delta$ is built from the estimated response
$\hat\eps(\cdot)$. Estimation error changes the closed-loop map itself;
this is where safety must be engineered, and it is the subject of L4.

Second, basin exit, the nonlinear channel: large probes can leave the
contraction neighbourhood. This is handled by the trust region
$|\dev_t| \le r$, a constraint the design problems of D4 already carry.

\paragraph{The perturbed-modulus lemma (L4).}
Let the corrected map be
\begin{equation*}
\Psi(h) \;=\; h + \eta\,\big[\,G(h) - \beta\,(h-\psi)\,\hat\eps(h)\,\big],
\end{equation*}
let $\hat\rho$ be its slope at the fixed point, and let $\rho$ be the
slope under the true response $\eps(\cdot)$, namely
$\rho = 1 - \eta\,\gpo$, the Newton-scaled contraction of the corrected
update \citep{reflex2026}.

\textbf{Lemma L4 (perturbed modulus).} With
$e(h) := \hat\eps(h) - \eps(h)$ the $C^1$ estimation error,
\begin{equation*}
\big|\hat\rho - \rho\big|
\;\le\;
\eta\,\beta\,\Big( \|e\|_\infty + |\hstar - \psi|\,\|e'\|_\infty \Big)
\;=:\;
\eta\,\beta\,L_{\mathrm{corr}}\,\|e\|_{C^1} .
\end{equation*}
The computation is short. The slope of the corrected map is
$\Psi'(h) = 1 + \eta\,\hat\Phi''(h)$, so the modulus perturbation is
$\eta$ times the curvature error, and
\begin{align*}
\hat\Phi''(h) - \Phi''(h)
&= -\beta\,\frac{d}{dh}\Big[\,(h-\psi)\,e(h)\,\Big] \\
&= -\beta\,\Big[\, e(h) + (h-\psi)\,e'(h) \,\Big].
\end{align*}
Taking absolute values and the supremum over the operating interval gives
the bound.

For the structural family the $C^1$ error is controlled by the parameter
error with explicit constants: with
$e = \eps(\cdot\,;\hat\theta) - \eps(\cdot\,;\theta)$,
\begin{equation*}
\|e\|_{C^1} \;\le\; L_\theta\,\|\hat\theta - \theta\| ,
\end{equation*}
where $L_\theta$ is computable from $(C_1, c, r)$. Consequently a confidence
region for $\theta$ translates directly into an interval for the
closed-loop modulus.

\paragraph{The pessimism rule and the safety proposition (P7.1).}
Maintain any anytime-valid confidence sequence $\Theta_t$ for the response
parameters (from common-random-number machinery as in \citet{reflex2026},
or a standard self-normalized bound; its width shrinks as information accumulates, per
D2).

\textbf{Rule.} At each step choose the pair (step size $\eta_t$, design)
so that the worst-case modulus over $\Theta_t$ respects the margin,
\begin{equation*}
\sup_{\theta \in \Theta_t}\; \big|\hat\rho(\theta)\big|
\;\le\; 1 - c_{\mathrm{margin}} ,
\end{equation*}
that is, pessimistic on stability, while the design objective of D4 is
evaluated optimistically. When $\Theta_t$ is too wide to certify the
margin, the rule freezes the correction and explores under the blind loop,
which is provably stable with modulus $m<1$ by R2. This is exactly the
anti-echo freeze of \texttt{structural\_response.py}, now derived rather
than engineered: this freeze is the pessimism gate.

\textbf{Proposition P7.1 (safety).} If the confidence sequence is valid
at level $1-\delta$ uniformly over time, then with probability at least
$1-\delta$ the realized modulus satisfies
$|\rho_t| \le 1 - c_{\mathrm{margin}}$ for all $t$ simultaneously. The
argument is one line: on the event that $\theta \in \Theta_t$ for every
$t$, which has probability at least $1-\delta$ by anytime-validity, the
rule's supremum dominates the realized modulus at every step. The
engineering reading: safety is inherited from the confidence sequence's
uniform validity; there is no union bound over steps and no per-step
failure accumulation.

\paragraph{Design regret (stated; journal deliverable).}
\textbf{Claim (to be proved in full).} The certainty-equivalent
re-solving scheme (the D-optimal objective from D4, the pessimism rule
above, and the trust region $r$) attains identification regret
$O(\sqrt{T})$ against the budget-matched oracle design. The
elliptical-potential argument of linear-bandit design theory supplies the
rate, and the freeze episodes contribute an additive $O(\log T)$ term:
each freeze ends when the confidence width halves, and the widths shrink
geometrically in the accumulated design energy. As a remark on scope,
this item is labeled partial for the workshop version: the algorithm and
the safety guarantee are proved, while the regret constant and the
freeze-episode accounting are the journal version's work.

\paragraph{Numerical verification (V6).}
Three checks, all in V6.
\begin{itemize}
\item Open-loop neutrality: injected bounded probe schedules never change
the measured homogeneous contraction rate, read off as the slope of the
deviation envelope, confirming R2.
\item The perturbed-modulus lemma: corrected loops were run with
deliberately corrupted $\hat\eps$, the $C^1$ error swept over a grid; in
every cell the measured fixed-point slope deviates from $1 - \eta\,\gpo$
by less than the bound of L4, and the bound is tight within a factor of
about two at small errors.
\item The pessimism rule: under a simulated valid confidence sequence, no
trajectory in 400 seeded runs violates the margin; under a deliberately
invalid (too-narrow) sequence, violations appear, so the validity
assumption of P7.1 is necessary, not an artifact of the proof.
\end{itemize}

\subsection*{D7: The anchoring crossover (T7)}

This subsection records the derivation behind T7: when does anchoring the
response estimate to the structural family (``structure'') beat a free-form
nonparametric estimator (``capacity'')? The headline of the derivation
record is that the crossover is horizon-dependent, not budget-only. The
formal statement is recorded as T7 in the theorem register. A status remark
carried over from the derivation record: the derivation is complete,
including the subtlety that changes the statement of the theorem (the
nonparametric bias floor is horizon-dependent), and both the exact
secant-bias constant and the two MSE formulas below are verified
numerically in V7.

\paragraph{The question.}
REFLEX's central empirical finding is that anchoring beats a free-form
estimator: the v3/v4 negative-then-positive result. When is that the right choice
\emph{in general}? The naive answer prices misspecification against the
exchange rate of T2 alone:
\begin{equation*}
  \text{anchor} \quad \text{iff} \quad
  \delta_{\mathrm{mis}}^2 \;<\; \frac{\gpo\,\sigma^2}{2B}.
\end{equation*}
Working the derivation shows the naive answer is incomplete: the
nonparametric alternative's bias is horizon-dependent, and the honest
crossover involves $(\delta_{\mathrm{mis}}, B, T)$ jointly.

\paragraph{The nonparametric side: local (secant) slope estimation.}
Estimate $\eps = -\tau'(\hstar)$ from symmetric probes at $\hstar \pm w$,
using $n$ probe pairs, under the value budget and the horizon. Writing
$S = 2 n w^2$ for the design energy, the two constraints read
\begin{equation*}
\begin{aligned}
  &\text{budget:} \quad \tfrac{1}{2}\,\gpo\,S \;\le\; B, \\
  &\text{horizon:} \quad 2n \;\le\; T
  \;\;\Longrightarrow\;\;
  w^2 \;=\; \frac{S}{2n} \;\ge\; \frac{S}{T},
\end{aligned}
\end{equation*}
and with the budget binding, $S = 2B/\gpo$, the horizon forces fat probes:
\begin{equation*}
  w^2 \;\ge\; \frac{2B}{\gpo\,T}.
\end{equation*}

Secant bias, with its exact leading constant. For the symmetric difference
quotient,
\begin{equation*}
  \frac{\tau(\hstar+w) - \tau(\hstar-w)}{2w}
  \;=\; \tau'(\hstar) \;+\; \frac{\tau'''(\hstar)}{6}\,w^2 \;+\; O(w^4):
\end{equation*}
the even-order terms cancel by symmetry, so the second derivative does not
enter at all. Symmetric probing is therefore already bias-optimal at its
order, a point worth one sentence in the paper. For the structural family
$\tau(h) = C_0 + C_1 e^{-ch}$ one has $\tau'''(h) = -c^3 C_1 e^{-ch}$, so
\begin{equation*}
  \mathrm{bias}(w)
  \;=\; -\frac{c^3 C_1 e^{-c\hstar}}{6}\,w^2
  \;=\; -\frac{c^2}{6}\,\eps\,w^2,
\end{equation*}
using $\eps = \eps(\hstar) = -\tau'(\hstar) = c\,C_1 e^{-c\hstar}$: a
curvature multiple of the slope itself, with the sign fixed by the family
(only the magnitude enters the MSE below).

MSE at the budget- and horizon-constrained optimum. Each probe pair yields
a difference quotient with variance $\sigma^2/(2w^2)$; averaging the $n$
pairs gives
\begin{equation*}
  \Var \;=\; \frac{\sigma^2}{2 n w^2} \;=\; \frac{\sigma^2}{S}
  \;=\; \frac{\gpo\,\sigma^2}{2B}
\end{equation*}
at the binding budget, independent of how $S$ splits into $n$ and $w$. The
bias, in contrast, requires small $w$, and the horizon caps how small:
$w^2 \ge 2B/(\gpo T)$. Hence, at the constrained optimum,
\begin{equation*}
  \mathrm{MSE}_{\mathrm{np}}
  \;=\; \frac{\gpo\,\sigma^2}{2B}
  \;+\; \Big(\frac{\tau'''(\hstar)}{6}\Big)^{\!2}
        \Big(\frac{2B}{\gpo\,T}\Big)^{\!2}.
\end{equation*}

The discovered subtlety, and the reason T7 carries the horizon: the bias
term \emph{falls} as $T$ grows at fixed budget. With more (cheaper,
smaller) probes the nonparametric estimator approaches the parametric
rate. Nonparametric identification is not uniformly worse; it is worse at
\emph{short horizons and fat probes}.

\paragraph{The parametric (anchored) side.}
For a correctly specified family,
$\mathrm{MSE}_{\mathrm{p}} = \kappa_p\,\gpo\,\sigma^2/(2B)$, where
$\kappa_p \ge 1$ is the design efficiency for extracting $\tau'(\hstar)$
through the family; c-optimal probing in the sense of T5(c), worked in D4,
makes $\kappa_p \approx 1$, and the constant is computable. Misspecified by
$\delta_{\mathrm{mis}}$, the $C^1$ distance from the true response to the
family as in D6, the anchored estimator picks up an irreducible squared
bias $\delta_{\mathrm{mis}}^2$:
\begin{equation*}
  \mathrm{MSE}_{\mathrm{anchor}}
  \;=\; \kappa_p\,\frac{\gpo\,\sigma^2}{2B} \;+\; \delta_{\mathrm{mis}}^2.
\end{equation*}

\paragraph{The crossover.}
Anchor iff $\mathrm{MSE}_{\mathrm{anchor}} < \mathrm{MSE}_{\mathrm{np}}$,
that is,
\begin{equation*}
  \delta_{\mathrm{mis}}^2 \;<\;
  \Big(\frac{\tau'''(\hstar)}{6}\Big)^{\!2}
  \Big(\frac{2B}{\gpo\,T}\Big)^{\!2}
  \;+\; (1 - \kappa_p)\,\frac{\gpo\,\sigma^2}{2B};
\end{equation*}
since $\kappa_p \ge 1$, the second term is a (small) penalty on anchoring
that vanishes under c-optimal probing through the family. The dominant
first term gives the memorable form stated in T7:
\begin{equation*}
  \text{anchor} \quad \text{iff} \quad
  \delta_{\mathrm{mis}} \;<\;
  \frac{|\tau'''(\hstar)|\,B}{3\,\gpo\,T}
  \quad \text{(to leading order)}.
\end{equation*}

\paragraph{Readings.}
(i) Anchoring wins at short horizons, tight budgets, and strongly curved
responses; the free-form route wins asymptotically in $T$ at fixed budget.
(ii) The formula retrodicts the v3 negative result \emph{quantitatively}:
the free-form loop failed at modest $T$ with the trust region forcing fat
effective probes, exactly the regime the formula assigns to anchoring.
(iii) It also predicts when the finding would reverse (long horizons,
generous probe counts), a regime no REFLEX experiment tests \citep{reflex2026}:
a falsifiable novel prediction. (iv) Delta over Lin--Zrnic (literature
review item G5): their plug-in analysis has neither the budget nor the
horizon; this formula prices both in the paper's single currency.

\paragraph{Numerical verification (V7).}
Three checks. (1) The secant-bias constant: the measured ratio
$(\text{secant} - \tau')/w^2$ converges to $\tau'''(\hstar)/6$ for the
exponential family (a deterministic check with exact function evaluations,
three values of $w$ extrapolated). (2) The $\mathrm{MSE}_{\mathrm{np}}$
formula: Monte-Carlo secant estimation at the constrained optimum matches
the two-term formula across a $(B, T)$ grid. (3) The crossover: simulated
anchored fits (correct family plus injected misspecification) against the
secant estimator show an empirical crossover $\delta_{\mathrm{mis}}^{*}$
tracking the formula within Monte-Carlo tolerance, including its $B/T$
scaling (halving $T$ doubles the crossover threshold).

\subsection*{D8/D9: The ROI of self-knowledge and the LQ and multi-bond instantiations (T8, P9.1, P9.2, P9.3)}

This subsection records the derivations behind Theorem T8 (the return on
investment of self-knowledge), the two consistency propositions P9.1 and
P9.2, and the multi-bond curvature instantiation P9.3. Status of the
underlying derivation document: the ROI closed forms are derived and
verified (V8); both consistency propositions are proved (one-line proofs,
verified); the multi-bond $\Gpo$ is derived exactly in the separable case,
with the coupled extension stated to first order. The formal statements
live in the theorem register and are quoted here only as needed for the
algebra.

\paragraph{Ingredients.}
Three constants, all computable from earlier derivations and the
primitives of the REFLEX market \citep{reflex2026}, enter the decision
problem:
\begin{equation*}
A = \tfrac{1}{2}\,\gpo\,(\hsp - \hpo)^2, \qquad
b = \tfrac{1}{2}\,\gpo\,\sigma^2, \qquad
a = \tfrac{1}{2}\,\gpo\,\kappa^2 .
\end{equation*}
Here $A$ is the echo-chamber gap, the per-step value of
running the corrected deployment $\hpo$ instead of the naive $\hsp$;
$b$ is the exchange rate of T2, under
which estimator accuracy $v$ costs $b/v$ in performative value
(derivation D2); and $a$ is the residual gap per unit of estimator
variance: an error in $\hat\eps$ of variance $v$ displaces the corrected
deployment by $\kappa^2 v$ in squared spread units, where
$\kappa = |\,d\hpo/d\eps\,|$ is the sensitivity of the
corrected optimum to the response slope, available in closed form from
the market model's fixed-point equation \citep{reflex2026}.

\paragraph{The decision problem.}
Explore once to accuracy $v$ (paying the exchange-rate price $b/v$ up
front), then run the corrected deployment forever, discounting the
perpetual stream at rate $\rho_{\mathrm{disc}}$. The corrected stream is
worth $A - a v$ per step, so the net present value is
\begin{equation*}
\mathrm{NPV}(v) \;=\; \frac{A - a v}{\rho_{\mathrm{disc}}} \;-\;
\frac{b}{v}, \qquad v \in (0,\infty).
\end{equation*}

\paragraph{Optimal accuracy and break-even patience (T8).}
Differentiating,
\begin{equation*}
\mathrm{NPV}'(v) \;=\; -\frac{a}{\rho_{\mathrm{disc}}} + \frac{b}{v^2}
\;=\; 0
\qquad\Longrightarrow\qquad
v^{*} \;=\; \sqrt{\frac{b\,\rho_{\mathrm{disc}}}{a}}
\;=\; \frac{\sigma}{\kappa}\,\sqrt{\rho_{\mathrm{disc}}},
\end{equation*}
where the second form follows by substituting $a$ and $b$: the common
factor $\tfrac{1}{2}\gpo$ cancels, so the optimal accuracy does not
involve the curvature at all. Since $\mathrm{NPV}''(v) = -2b/v^3 < 0$ on
$(0,\infty)$, the objective is concave there and $v^{*}$ is the global
maximizer. At the optimum the two $v$-dependent terms are equal,
\begin{equation*}
\frac{a\,v^{*}}{\rho_{\mathrm{disc}}}
= \sqrt{\frac{a\,b}{\rho_{\mathrm{disc}}}}
= \frac{b}{v^{*}},
\end{equation*}
so the optimal value is
\begin{equation*}
\mathrm{NPV}(v^{*}) \;=\; \frac{A}{\rho_{\mathrm{disc}}}
\;-\; 2\,\sqrt{\frac{a\,b}{\rho_{\mathrm{disc}}}} .
\end{equation*}
Exploring at all is worthwhile iff $\mathrm{NPV}(v^{*}) > 0$, which
rearranges to
\begin{equation*}
A \;>\; 2\,\sqrt{a\,b\,\rho_{\mathrm{disc}}}
\qquad\Longleftrightarrow\qquad
\rho_{\mathrm{disc}} \;<\; \rho^{*} \;:=\; \frac{A^2}{4\,a\,b}
\;=\; \frac{(\hsp - \hpo)^4}{4\,\kappa^2\,\sigma^2}.
\end{equation*}
The last equality is the substitution
$A^2/(4ab) = \big(\tfrac{1}{2}\gpo\big)^2 (\hsp-\hpo)^4 \,/\,
\big(4 \cdot \tfrac{1}{2}\gpo \kappa^2 \cdot \tfrac{1}{2}\gpo
\sigma^2\big)$, in which $\gpo$ cancels completely.

Three readings of the closed forms. (i) The break-even patience
$\rho^{*} = (\hsp-\hpo)^4/(4\kappa^2\sigma^2)$ involves only the
deployment gap, the correction's sensitivity, and the response noise:
the curvature drops out, a clean and slightly surprising
closed form (it cancels between the gap's value and the exploration
price). (ii) Everything on the right-hand side is computable from the
primitives of REFLEX theory before exploring, which is the point of the
result as a decision rule. (iii) $v^{*}$ is the right amount of
self-knowledge: more accuracy than $v^{*}$ is vanity, less is blindness.

\paragraph{Frontier consistency (P9.1).}
For any deterministic design $\dev_t$ with energy
$S_T = \sum_{t\le T}\dev_t^2$, the pathwise exchange-rate product is
\begin{equation*}
\frac{\sigma^2}{S_T} \times \tfrac{1}{2}\,\gpo\, S_T
\;=\; \tfrac{1}{2}\,\gpo\,\sigma^2 :
\end{equation*}
the $S_T$ cancels, so nothing about the schedule matters. In particular
the $\dev_t \sim c/\sqrt{t}$ schedules with $S_T = c^2 \log T$, the
Lai--Robbins cost-$O(\log T)$ regime of adaptive design, satisfy the
exchange rate of T2 with equality: their celebrated cheapness buys
proportionally little information, and they sit exactly on the frontier.

\paragraph{Isotropy contrast (P9.2).}
In the $d$-dimensional problem, the ratio of the isotropic design's risk
to the A-optimal design's risk at matched budget is the
curvature-dispersion factor of D4 and C5.1,
\begin{equation*}
F \;=\; \frac{d\,\tr(\Gpo)}{\big(\tr \Gpo^{1/2}\big)^{2}} \;\ge\; 1,
\end{equation*}
with $F = 1$ iff $\Gpo$ is a multiple of the identity. That is exactly
the LQR-like isotropic case in which naive isotropic exploration is
optimal (Simchowitz--Foster 2020); generically $F > 1$ and grows with
the spread of the curvature spectrum. The performative setting is
generically anisotropic (the bond-level curvatures $\gamma_a$ span
decades across the calibrated universe), so naive exploration is
generically suboptimal: this is the contrast between the two regimes
stated in the main text.

\paragraph{Separable multi-bond curvature (P9.3, exact).}
In the multi-bond market with per-bond decisions $h_a$ and separable
flows, the response Jacobian is diagonal,
$E = \operatorname{diag}(\eps_a)$, and the performative objective
separates across bonds. Applying the scalar curvature computation
\citep{reflex2026} bond by bond gives, exactly,
\begin{equation*}
(\Gpo)_{aa} \;=\; \gamma_a \;+\; \beta\,\eps_a\,
\big( 2 + c_t\,\psi_a - c_t\,h_a \big),
\qquad (\Gpo)_{ab} = 0 \;\; (a \ne b).
\end{equation*}
All quantities on the right are the calibrated per-bond constants of
the REFLEX market \citep{reflex2026}. Consequently the
dispersion factor $F$ of P9.2 is directly computable on the calibrated
universe (170 CUSIPs with at least twelve months of returns; the
REALDATA record); that computed value is the number quoted in the
paper's experiments.

\paragraph{Factor-coupled case (construction, first order).}
With factor-coupled responses, $E$ diagonal plus low rank as in REFLEX
theory 1.5, the same differentiation at the operating point gives
\begin{equation*}
\begin{aligned}
\Gpo \;&=\; \Gamma \;+\; \beta\,\Big[\, 2\,E_{\mathrm{sym}}
+ c_t\big( \operatorname{diag}(\psi)\,E_{\mathrm{sym}}
- H\,E_{\mathrm{sym}} \big) \Big]
\;+\; O\!\big(\|E_{\mathrm{offdiag}}\|^2\big), \\[2pt]
E_{\mathrm{sym}} \;&=\; \tfrac{1}{2}\,(E + E^{\top}), \qquad
H = \operatorname{diag}(h_a),
\end{aligned}
\end{equation*}
which is exact when $E$ is symmetric and is evaluated at the operating
point. The low-rank structure of $E$ is inherited by the correction, so
$\Gpo^{-1/2}$ and $\tr \Gpo^{1/2}$ remain computable in $O(d k^2)$ via
the same Woodbury and eigen-split machinery already used in REFLEX
theory 1.5: the instantiation adds no new computational burden.

Caveat (honesty convention). The paper states the separable case as an
exact result (P9.3) and the coupled case only as a first-order
construction carrying its explicit $O(\|E_{\mathrm{offdiag}}\|^2)$ error
term, matching the reporting conventions of \citet{reflex2026}.

\paragraph{Numerical verification (V8).}
Five checks. (1) $v^{*} = \sqrt{b\,\rho_{\mathrm{disc}}/a}$ and the
$\mathrm{NPV}(v^{*})$ formula match a direct numerical scan over $v$ on
three parameter sets. (2) Root-finding
$\mathrm{NPV}(v^{*})(\rho_{\mathrm{disc}}) = 0$ in $\rho_{\mathrm{disc}}$
recovers $\rho^{*} = (\hsp-\hpo)^4/(4\kappa^2\sigma^2)$ to $10^{-10}$,
confirming the $\gpo$ cancellation. (3) The frontier proposition:
$c/\sqrt{t}$ schedules measure exactly on-frontier. (4) The dispersion
factor on a synthetic curvature spectrum spanning two decades: $F$
matches the measured isotropic-to-optimal risk ratio (cross-checked
against V4). (5) The separable $\Gpo$: a finite-difference Hessian of
the separable objective matches the per-bond formula entrywise.
